\documentclass[11pt]{article}
\usepackage[margin=0.86in]{geometry}
\usepackage{amsmath,amssymb,amsthm,graphicx,microtype,enumitem}
\usepackage[hidelinks]{hyperref}

\newtheorem{theorem}{Theorem}
\newtheorem{lemma}[theorem]{Lemma}
\theoremstyle{remark}
\newtheorem{remark}[theorem]{Remark}
\newcommand{\T}{\mathbb T}
\newcommand{\Z}{\mathbb Z}
\newcommand{\R}{\mathbb R}

\title{Eventual ultraboundedness has a sharp nonlocal threshold}
\author{Candidate sharp partial obstruction for arXiv:2002.09221}
\date{13 August 2026}

\begin{document}
\maketitle

\begin{center}
\fbox{\parbox{0.92\textwidth}{\textbf{Status.}
Candidate sharp partial obstruction, likely valid.  We give an exact
counterexample to the broad reversible-Markov implication suggested by the
source: a kernel can be bounded at one positive time without being bounded at
all positive times.  The example is nonlocal, so the source's question for
smooth Langevin diffusions on Euclidean space remains open here.}}
\end{center}

\section{Source and exact signal}

The source is P. Cattiaux, M. Fathi, and A. Guillin, \emph{Self-improvement
of the Bakry--Emery criterion for Poincar\'e inequalities and Wasserstein
contraction using variable curvature bounds}, arXiv:2002.09221 \cite{CFG}.
On page 26, in the setting of the reversible diffusion generated by
$L=\Delta-\nabla V\cdot\nabla$ on $\R^n$, the authors write:

\begin{figure}[ht]
\centering
\includegraphics[width=0.98\textwidth]{figures/open_problem_crop.png}
\caption{The exact open signal on page 26 of the source PDF.}
\end{figure}

\noindent\textbf{Transcription.}
``Actually, we do not know examples where the semi-group is not ultra
bounded, but $r(2\eta,x,x)$ is bounded in $x$.''  Here $r(t,x,y)$ is the
transition density with respect to the invariant probability $\mu$;
``ultra bounded'' means
$\sup_{x,y}r(t,x,y)<\infty$ for every $t>0$.

\section{The exact reduction}

\begin{lemma}[one-time diagonal bound equals eventual ultraboundedness]
Let $(P_t)_{t\geq0}$ be a symmetric Markov semigroup with densities $r_t$
with respect to a probability measure $\mu$.  If
\[
 M_T:=\sup_x r_T(x,x)<\infty
\]
for one $T>0$, then $\sup_{x,y}r_s(x,y)\leq M_T$ for every $s\geq T$.
\end{lemma}

\begin{proof}
Symmetry and the semigroup identity give
\[
 r_T(x,y)=\int r_{T/2}(x,z)r_{T/2}(y,z)\,d\mu(z).
\]
Cauchy--Schwarz therefore yields
$r_T(x,y)^2\leq r_T(x,x)r_T(y,y)\leq M_T^2$.  For $s\geq T$,
\[
 r_s(x,y)=\int r_T(x,z)r_{s-T}(z,y)\,d\mu(z)\leq M_T,
\]
because the second factor integrates to one.
\end{proof}

Thus the source is asking whether \emph{eventual} ultraboundedness can occur
without \emph{immediate} ultraboundedness in its local diffusion class.

\section{Proof intuition}

Subordination slows the high-frequency damping of a heat semigroup.  On the
torus, Brownian eigenvalues grow like $|k|^2$.  Replacing them by
$\log(1+|k|^2)$ changes exponential damping into the polynomial multiplier
$(1+|k|^2)^{-t}$.  This multiplier is summable only after the Sobolev
threshold $2t>d$.  Gamma subordination preserves positivity, constants,
symmetry, and the Markov property, so every probabilistic feature except
locality survives.  A periodized-Gaussian lower bound proves that the failure
at and below the threshold is genuine kernel blow-up, not merely failure of
an absolute-convergence test.

\section{The sharp threshold example}

Let $\T^d=\R^d/\Z^d$ have normalized Haar measure, let
$A=-\Delta\geq0$, and define
\[
 P_t=e^{-t\log(I+A)}=(I-\Delta)^{-t},\qquad t\geq0.
\]

\begin{theorem}[exact eventual-ultraboundedness threshold]\label{thm:main}
The family $(P_t)_{t\geq0}$ is a conservative symmetric Markov semigroup on
$\T^d$.  Its transition kernel is bounded if and only if
\[
 t>\frac d2.
\]
Consequently, for every $\eta>d/4$ one has
$\sup_x r_{2\eta}(x,x)<\infty$, but the semigroup is not ultrabounded in the
all-positive-times sense used in arXiv:2002.09221.
\end{theorem}

\begin{proof}
For $t>0$ and $\lambda\geq0$, the Gamma integral gives
\[
 (1+\lambda)^{-t}=\frac1{\Gamma(t)}\int_0^\infty
 s^{t-1}e^{-s}e^{-s\lambda}\,ds.                         \tag{4.1}
\]
Hence
\[
 P_t=\frac1{\Gamma(t)}\int_0^\infty s^{t-1}e^{-s}e^{s\Delta}\,ds.
                                                                    \tag{4.2}
\]
This is a probability mixture of the Brownian heat operators.  It is
therefore positivity preserving and fixes constants.  Self-adjointness is
clear, and the spectral multipliers show $P_tP_u=P_{t+u}$.

The Fourier kernel is
\[
 K_t(z)=\sum_{k\in\Z^d}(1+4\pi^2|k|^2)^{-t}e^{2\pi i k\cdot z}.
                                                                    \tag{4.3}
\]
If $t>d/2$, the coefficients belong to $\ell^1(\Z^d)$, so (4.3)
converges absolutely and uniformly to a bounded continuous kernel.

It remains to prove necessity.  Let $h_s$ denote the torus heat kernel.  On
the representative cube $[-1/2,1/2]^d$,
\[
 h_s(z)=(4\pi s)^{-d/2}\sum_{m\in\Z^d}
          \exp\!\left(-\frac{|z+m|^2}{4s}\right).
\]
Combining this with (4.2), retaining only $m=0$, and taking
$0<|z|\leq1/4$, we obtain
\begin{align*}
 K_t(z)
 &\geq \frac1{\Gamma(t)}\int_{|z|^2}^{1}
 s^{t-1}e^{-s}(4\pi s)^{-d/2}
 \exp\!\left(-\frac{|z|^2}{4s}\right)ds\\
 &\geq c_{d,t}\int_{|z|^2}^{1}s^{t-1-d/2}\,ds.             \tag{4.4}
\end{align*}
When $t<d/2$, the last expression grows like
$|z|^{2t-d}$; when $t=d/2$, it grows like $\log(1/|z|)$.
Thus $K_t$ is unbounded on every neighborhood of zero for all $t\leq d/2$.
This proves the exact threshold.  Choosing $2\eta>d/2$ gives the claimed
diagonal bound, while any $0<t\leq d/2$ disproves all-time
ultraboundedness.
\end{proof}

\section{What the example settles---and what it does not}

Theorem~\ref{thm:main} shows that no abstract argument using only symmetry,
reversibility, the Markov property, compact state space, or a probability
invariant measure can propagate a one-time kernel bound backwards to every
positive time.  General eventual-ultracontractivity mechanisms under
subordination are established in \cite{BCS,GM}; the exact torus threshold is
included here with a self-contained proof.  Novelty of that specialization
is not asserted.

The missing source hypothesis is locality.  The generator
$-\log(I-\Delta)$ is nonlocal (the process is Gamma-subordinated Brownian
motion), whereas arXiv:2002.09221 works with the second-order diffusion
$\Delta-\nabla V\cdot\nabla$ on $\R^n$.  Therefore this packet is a sharp
partial obstruction, not a full counterexample to the source-context
question.

Eight focused upgrades were attempted.  Besides the torus construction, an
exact reversible star chain gives the same threshold phenomenon.  A
ground-state-transform search found nonlocal but no decisive local examples.
A one-dimensional hitting-time route encountered an explicit correction in
the authors' later paper \cite{CG}: the earlier bounded-Lyapunov implication
requires extra assumptions.  Standard quadratic and superquadratic
potentials exhibit an all-times dichotomy.  Finally, narrow remote plateaux
can reproduce the necessary local lower bounds, but a uniform large-time
heat-kernel upper bound across all connecting regions remains unproved.

\section{Verification, novelty bounds, and review recommendation}

The proof has two independent halves: $\ell^1$ Fourier summability above the
threshold and a positive heat-kernel lower bound below it.  The latter also
checks the endpoint $t=d/2$, where the singularity is logarithmic.  The
included verifier report records these checks and the scope boundary.

The run's four lightweight indexes were searched for arXiv:2002.09221,
eventual ultracontractivity, Gamma subordination, and the core kernel phrase;
no duplicate packet was found.  Bounded external searches through
13 August 2026 included the exact quoted sentence, eventual/asymptotic
ultracontractivity, classical Schr\"odinger variants, and bounded-Lyapunov
criteria.  They found the general subordinated-semigroup literature
\cite{BCS,GM} and the correction \cite{CG}, but no resolution of the smooth
finite-dimensional Langevin case.

Human review should check (i) that the intended scope of the source sentence
indeed retains its surrounding Langevin assumptions, (ii) the endpoint in
(4.4), and (iii) the classification as a partial obstruction.  Promotion to
a full answer requires either a smooth local diffusion with this threshold
behavior or a proof that eventual ultraboundedness implies immediate
ultraboundedness for all such diffusions.

\begin{thebibliography}{9}
\bibitem{CFG}
P. Cattiaux, M. Fathi, and A. Guillin,
\emph{Self-improvement of the Bakry--Emery criterion for Poincar\'e
inequalities and Wasserstein contraction using variable curvature bounds},
arXiv:2002.09221.

\bibitem{BCS}
A. Bendikov, T. Coulhon, and L. Saloff-Coste,
\emph{Ultracontractivity and embedding into $L^\infty$},
Math. Ann. 337 (2007), 817--853.

\bibitem{GM}
I. Gentil and P. Maheux,
\emph{Super-Poincar\'e and Nash-type inequalities for subordinated
semigroups}, arXiv:1105.3095.

\bibitem{CG}
P. Cattiaux and A. Guillin,
\emph{Hitting times, functional inequalities, Lyapunov conditions and
uniform ergodicity}, J. Funct. Anal. 272 (2017), 2361--2391;
arXiv:1604.06336.
\end{thebibliography}

\end{document}
